\chapter{Information-Theoretic Cryptography}

In this chapter the adversary is \textbf{computationally unbounded}: she has infinite time and memory. Every result holds without any assumption, which is why this kind of security is called \textbf{information-theoretic} (or \textit{unconditional}). We will see that this strong guarantee is possible, both for confidentiality and for integrity, but that it always comes with a price: keys must be very long.

% ============================================================
%  Symmetric-key encryption, perfect secrecy, One-Time Pad, Shannon
% ============================================================

\section{Symmetric-key encryption}\label{sec:ske}

Let's fix the scene that we will use for the whole chapter. \textbf{Alice} wants to send a message to \textbf{Bob} over a public channel. \textbf{Eve}, the eavesdropper, reads everything that travels on that channel. Alice and Bob have one advantage over Eve: they share a \textbf{secret key}, chosen uniformly at random, that Eve does not know. Since the \textbf{same key is used both to encrypt and to decrypt}, this setting is called \textbf{symmetric} (or private-key).

\begin{figure}[H]
  \centering
  \begin{tikzpicture}[node distance=1.2cm, box/.style={draw, rounded corners=2pt, minimum height=7mm, inner sep=4pt}]
    \node (A) {Alice};
    \node[box, right=1.2cm of A] (E) {$\Enc(k, m)$};
    \coordinate[right=1.6cm of E] (mid);
    \node[box, right=1.6cm of mid] (D) {$\Dec(k, c)$};
    \node[right=1.2cm of D] (B) {Bob};
    \node[below=1.1cm of mid] (Eve) {Eve};
    \node[below=0.9cm of E] (k1) {$k$};
    \node[below=0.9cm of D] (k2) {$k$};
    \draw[msg] (A) -- node[above] {$m$} (E);
    \draw[thick] (E) -- node[above] {$c$} (mid);
    \draw[msg] (mid) -- (D);
    \draw[msg, dashed] (mid) -- (Eve);
    \draw[msg] (D) -- node[above] {$m$} (B);
    \draw[msg] (k1) -- (E);
    \draw[msg] (k2) -- (D);
  \end{tikzpicture}
  \caption{Symmetric-key encryption. Eve sees the ciphertext $c$ but not the shared key $k$.}
\end{figure}

The objects we work with:
\begin{itemize}
  \item \textbf{Key space} $\Kspace$ $\to$ the set of possible keys; the key is sampled uniformly, $K \sample \Kspace$;
  \item \textbf{Message space} $\Mspace$ $\to$ the set of possible plaintexts;
  \item \textbf{Ciphertext space} $\Cspace$ $\to$ the set of possible ciphertexts.
\end{itemize}

\begin{frameddefn}[label={def:ske}]{Symmetric-key encryption scheme}
  A symmetric-key encryption scheme is a pair $\Pi = (\Enc, \Dec)$ with
  \[
    \Enc : \Kspace \times \Mspace \to \Cspace \qquad \Dec : \Kspace \times \Cspace \to \Mspace
  \]
  that satisfies \textbf{correctness}: decrypting with the right key always gives back the original message,
  \[
    \forall m \in \Mspace,\ \forall k \in \Kspace : \quad \Dec(k, \Enc(k, m)) = m
  \]
\end{frameddefn}

Correctness is only the minimum requirement: it says that the scheme \textit{works}, not that it is \textit{secure}. The scheme $\Enc(k, m) = m$ is perfectly correct and completely useless.

\subsection{Kerckhoffs's principle}

Where should the security of a scheme come from? Kerckhoffs's answer (1883), still valid today:

\begin{framedobs}{Kerckhoffs's principle}
  The security of a scheme must depend \textbf{only on the secrecy of the key}, never on the secrecy of the algorithm. Shannon summarized it as \textit{``the enemy knows the system''}.
\end{framedobs}

Why is this the right choice?
\begin{itemize}
  \item \textbf{Algorithms leak} $\to$ software gets reverse-engineered, devices get stolen, employees change jobs. A key is short and easy to replace; an algorithm is not.
  \item \textbf{Public algorithms get tested} $\to$ a public algorithm can be analysed by the whole research community, which is the best way to find weaknesses before an attacker does.
\end{itemize}
From now on, Eve always knows exactly how $\Enc$ and $\Dec$ work. The only thing she does not know is the key.

\section{Perfect secrecy}\label{sec:perfect-secrecy}

We now want to formalize the idea that ``Eve learns nothing from the ciphertext''. Think about what Eve knows \textbf{before} seeing anything: she is not completely ignorant. She may know that the message is probably ``attack at dawn'' and unlikely to be ``attack at noon''. This is her \textbf{a priori} knowledge, and we model it as a probability distribution over the messages.

Then Eve sees the ciphertext $c$ and updates her knowledge: this is the \textbf{a posteriori} knowledge. Shannon's idea (1949) is simple: \textbf{the scheme is perfectly secret if seeing the ciphertext does not change Eve's knowledge at all.}

The formal setup:
\begin{itemize}
  \item \textbf{Message} $\to$ a random variable $M$ over $\Mspace$ with an \textbf{arbitrary} distribution (what Eve knows a priori);
  \item \textbf{Key} $\to$ $K \sample \Kspace$, uniform and \textbf{independent of} $M$ (the key is chosen regardless of the message);
  \item \textbf{Ciphertext} $\to$ the random variable $C = \Enc(K, M)$ that we obtain from them.
\end{itemize}

\begin{frameddefn}[label={def:perfect-secrecy}]{Perfect secrecy}
  $\Pi$ is \textbf{perfectly secret} if, \textbf{for every distribution of $M$}, for every $m \in \Mspace$ and every $c \in \Cspace$ with $\Pr[C = c] > 0$:
  \[
    \Pr[M = m \mid C = c] = \Pr[M = m]
  \]
  In words: the a posteriori probability of every message is equal to its a priori probability.
\end{frameddefn}

Two details of this definition are essential:
\begin{itemize}
  \item \textbf{``For every distribution of $M$''} $\to$ security is a property of the \textbf{scheme}, and it must hold whatever Eve already knows about the message. We will see in \Cref{sec:equivalent} that without this quantifier the definition would be meaningless.
  \item \textbf{The adversary is unbounded} $\to$ nothing in the definition talks about time or resources: Eve can compute anything she wants.
\end{itemize}

\paragraph{Notation used in class.} The joint probability is written with the symbol $\wedge$ (``and''): $\Pr[M = m \wedge C = c]$ is the probability that the message is $m$ \textbf{and} the ciphertext is $c$. It is the same as $\Pr[M = m, C = c]$. With this notation, the left side of the definition is computed as
\[
  \Pr[M = m \mid C = c] = \frac{\Pr[M = m \wedge C = c]}{\Pr[C = c]}
\]
Knowing that the ciphertext is $c$ \textbf{restricts the world} to the outcomes where the ciphertext is $c$, and among those we measure in which fraction the message is $m$.

\begin{framedex}{Perfect secrecy with numbers}
  A general sends one bit with the One-Time Pad: ``attack'' (0) with probability $0.9$, ``retreat'' (1) with probability $0.1$. The key is 0 or 1, each with probability $0.5$. Eve sees the ciphertext 0. How likely is ``attack''?
  \begin{itemize}
    \item \textbf{Numerator} $\to$ ``attack'' \textbf{and} ciphertext 0 needs key 0: $\Pr[M = \text{attack} \wedge C = 0] = 0.9 \cdot 0.5 = 0.45$;
    \item \textbf{Denominator} $\to$ the ciphertext is 0 in two cases, ``attack'' with key 0 ($0.45$) or ``retreat'' with key 1 ($0.1 \cdot 0.5 = 0.05$), so $\Pr[C = 0] = 0.5$.
  \end{itemize}
  \[
    \Pr[M = \text{attack} \mid C = 0] = \frac{0.45}{0.5} = 0.9
  \]
  Before seeing the ciphertext Eve believed ``attack'' with probability $0.9$, and after seeing it she still believes $0.9$: \textbf{she learned nothing}. With a scheme that does not encrypt at all (ciphertext equal to the message) she would get $\frac{0.9}{0.9} = 1$: she would learn everything.
\end{framedex}

\section{Three equivalent definitions of perfect secrecy}\label{sec:equivalent}

The definition above describes what we want, but it is not very handy when we have to check a concrete scheme. The following lemma gives two equivalent ways to say the same thing, and the third one will be our main tool.

\begin{framedlem}[label={lem:equivalent}]{Equivalent definitions of perfect secrecy}
  For every symmetric-key encryption scheme $\Pi$, the following statements are equivalent:
  \begin{enumerate}[label=(\roman*)]
    \item \textbf{Perfect secrecy}: $\Pr[M = m \mid C = c] = \Pr[M = m]$ for every distribution of $M$;
    \item \textbf{Independence}: $M$ and $C$ are independent for every distribution of $M$, that is,
    \[
      \Pr[M = m \wedge C = c] = \Pr[M = m] \cdot \Pr[C = c]
    \]
    \item \textbf{Same ciphertext distribution}: for all $m, m' \in \Mspace$ and all $c \in \Cspace$,
    \[
      \Pr_K[\Enc(K, m) = c] = \Pr_K[\Enc(K, m') = c]
    \]
  \end{enumerate}
\end{framedlem}

How to read the three statements:

\begin{table}[H]
  \centering
  \small
  \begin{tabularx}{\textwidth}{@{}lXX@{}}
    \toprule
    \textbf{Statement} & \textbf{What it talks about} & \textbf{When it is useful} \\
    \midrule
    (i) Perfect secrecy & Eve's knowledge, a priori versus a posteriori & the \textit{intuitive} definition: what we actually want \\
    (ii) Independence & the joint distribution of $M$ and $C$ & a bridge between (i) and (iii); in information theory it means that the mutual information is zero, $I(M;C) = 0$ \\
    (iii) Same ciphertext distribution & only the randomness of the \textbf{key}: no distribution of $M$ is involved & the \textit{operational} one: the easiest to check on a concrete scheme \\
    \bottomrule
  \end{tabularx}
\end{table}

Statement (iii) deserves one more sentence: \textbf{the distribution of the ciphertext does not depend on the message}. Whatever message Alice encrypts, the ciphertext ``looks the same'' to Eve, so observing it cannot tell her which message was sent.

We prove the equivalence with a cycle of implications: (i) $\Rightarrow$ (ii) $\Rightarrow$ (iii) $\Rightarrow$ (i).

\begin{proof}[Proof of (i) $\Rightarrow$ (ii)]
  \textbf{Goal:} starting from perfect secrecy, obtain $\Pr[M = m \wedge C = c] = \Pr[M = m] \cdot \Pr[C = c]$.

  We start from hypothesis (i) and expand the conditional probability with its definition:
  \[
    \Pr[M = m] \overset{\text{(i)}}{=} \Pr[M = m \mid C = c] = \frac{\Pr[M = m \wedge C = c]}{\Pr[C = c]}
  \]
  The first equality is hypothesis (i); the second is the definition of conditional probability, which can be used because $\Pr[C = c] > 0$. Multiplying both sides by $\Pr[C = c]$:
  \[
    \Pr[M = m \wedge C = c] = \Pr[M = m] \cdot \Pr[C = c]
  \]
  which is exactly independence. If $\Pr[C = c] = 0$, both sides are $0$ and the equality holds anyway.
\end{proof}

\begin{proof}[Proof of (ii) $\Rightarrow$ (iii)]
  \textbf{Goal:} show that $\Pr[\Enc(K, m) = c] = \Pr[\Enc(K, m') = c]$ for any two messages. \textbf{Idea:} show that \textbf{both} are equal to the same quantity, $\Pr[C = c]$. We start from the left side and transform it step by step:
  \[
    \Pr[\Enc(K, m) = c] \overset{(a)}{=} \Pr[\Enc(K, M) = c \mid M = m] \overset{(b)}{=} \Pr[C = c \mid M = m] \overset{(c)}{=} \Pr[C = c]
  \]
  \begin{itemize}
    \item \textbf{(a)} ``Encrypting the \textit{fixed} message $m$'' is the same as ``encrypting the \textit{random} message $M$, knowing that it turned out to be $m$''. This works because the key $K$ is independent of $M$: knowing that $M = m$ tells us nothing about the key.
    \item \textbf{(b)} By definition $C = \Enc(K, M)$, so we are just renaming.
    \item \textbf{(c)} Hypothesis (ii): $M$ and $C$ are independent, so knowing that $M = m$ does not change the probability that $C = c$.
  \end{itemize}
  By the same steps with $m'$, $\Pr[\Enc(K, m') = c] = \Pr[C = c]$. Both probabilities are equal to $\Pr[C = c]$, so they are equal to each other.

  A technical detail: to condition on $M = m$ we need $\Pr[M = m] > 0$. Since (ii) holds \textbf{for every distribution}, we can choose $M$ uniform, so that every message has positive probability.
\end{proof}

\begin{proof}[Proof of (iii) $\Rightarrow$ (i)]
  \textbf{Goal:} from ``the ciphertext has the same distribution for every message'', obtain perfect secrecy. Fix any distribution of $M$, and $m, c$ with $\Pr[C = c] > 0$. The proof has two steps:
  \begin{enumerate}
    \item show that $\Pr[C = c] = \Pr[C = c \mid M = m]$: knowing the message does not change the probability of the ciphertext;
    \item use Bayes' rule to ``reverse'' this equality and obtain $\Pr[M = m] = \Pr[M = m \mid C = c]$.
  \end{enumerate}

  \textbf{Step 1.} We start from $\Pr[C = c]$ and transform it step by step:
  \[
    \begin{aligned}
      \Pr[C = c] &\overset{(a)}{=} \sum_{m'} \Pr[C = c \wedge M = m'] \\
      &\overset{(b)}{=} \sum_{m'} \Pr[C = c \mid M = m'] \cdot \Pr[M = m'] \\
      &\overset{(c)}{=} \sum_{m'} \Pr[\Enc(K, M) = c \mid M = m'] \cdot \Pr[M = m'] \\
      &\overset{(d)}{=} \sum_{m'} \Pr[\Enc(K, m') = c] \cdot \Pr[M = m'] \\
      &\overset{(e)}{=} \sum_{m'} \Pr[\Enc(K, m) = c] \cdot \Pr[M = m'] \\
      &\overset{(f)}{=} \Pr[\Enc(K, m) = c] \cdot \underbrace{\sum_{m'} \Pr[M = m']}_{=\,1} \\
      &\overset{(g)}{=} \Pr[\Enc(K, M) = c \mid M = m] = \Pr[C = c \mid M = m]
    \end{aligned}
  \]
  \begin{itemize}
    \item \textbf{(a)} Law of total probability: we split the event $C = c$ according to which message $m'$ was sent.
    \item \textbf{(b)} Definition of conditional probability, in the form $\Pr[A \wedge B] = \Pr[A \mid B] \cdot \Pr[B]$.
    \item \textbf{(c)} We write $C$ as what it is, $\Enc(K, M)$.
    \item \textbf{(d)} The same trick as step (a) of the previous proof: knowing that $M = m'$, encrypting $M$ is the same as encrypting the fixed message $m'$ (the key is independent of the message).
    \item \textbf{(e)} \textbf{Here we use hypothesis (iii)}: $\Pr[\Enc(K, m') = c]$ is the same for every message, so we can replace $m'$ with our fixed $m$.
    \item \textbf{(f)} $\Pr[\Enc(K, m) = c]$ no longer depends on $m'$, so it comes out of the sum. What is left is the sum of the probabilities of all messages, which is $1$.
    \item \textbf{(g)} The trick of (d) and (c) backwards: we go back to $\Pr[C = c \mid M = m]$.
  \end{itemize}

  \textbf{Step 2 (Bayes).} The probability of ``both events'' can be written in two ways, conditioning on one event or on the other:
  \[
    \Pr[M = m \mid C = c] \cdot \Pr[C = c] = \Pr[M = m \wedge C = c] = \Pr[C = c \mid M = m] \cdot \Pr[M = m]
  \]
  We solve for $\Pr[M = m]$ by dividing by $\Pr[C = c \mid M = m]$, which by Step 1 is equal to $\Pr[C = c] > 0$:
  \[
    \begin{aligned}
      \Pr[M = m] &= \frac{\Pr[M = m \mid C = c] \cdot \Pr[C = c]}{\Pr[C = c \mid M = m]} \\
      &= \frac{\Pr[M = m \mid C = c] \cdot \cancel{\Pr[C = c]}}{\cancel{\Pr[C = c]}} = \Pr[M = m \mid C = c]
    \end{aligned}
  \]
  This is perfect secrecy.
\end{proof}

\begin{framedobs}{Why ``for every distribution'' is essential}
  Suppose we dropped the quantifier and asked the definition to hold for \textbf{one single} distribution, for example the constant one $\Pr[M = m_0] = 1$ (Eve already knows the message). Then:
  \begin{itemize}
    \item statements (i) and (ii) would be true for \textbf{any} scheme, even $\Enc(k, m) = m$, which does not encrypt at all;
    \item statement (iii) would be false for that scheme.
  \end{itemize}
  So the lemma would be false. The step (ii) $\Rightarrow$ (iii) works only because we are free to choose the uniform distribution.
\end{framedobs}

\section{The One-Time Pad}\label{sec:otp}

We now have a definition and a handy way to check it. Does any scheme actually achieve perfect secrecy? Yes, and a very simple one: the \textbf{One-Time Pad} (Vernam, 1917). The idea is to hide every bit of the message by combining it with a \textbf{fresh, uniformly random key bit}, using the XOR operation (bitwise exclusive OR: $0$ if the two bits are equal, $1$ if they are different).

\begin{frameddefn}[label={def:otp}]{One-Time Pad}
  $\Mspace = \Kspace = \Cspace = \{0,1\}^n$, with the key sampled uniformly, $K \sample \{0,1\}^n$:
  \[
    \Enc(k, m) = k \oplus m \qquad \Dec(k, c) = k \oplus c
  \]
\end{frameddefn}

\textbf{Correctness.} We only need two properties of XOR: it is associative, and anything XOR-ed with itself gives zero ($k \oplus k = 0^n$):
\[
  \Dec(k, \Enc(k, m)) = k \oplus (k \oplus m) = (k \oplus k) \oplus m = 0^n \oplus m = m
\]

\begin{framedex}{One-Time Pad with four bits}
  Take $m = 1011$ and $k = 0110$. Then $c = 1101$, and decrypting gives back $c \oplus k = 1101 \oplus 0110 = 1011 = m$.

  Now look at it from Eve's side. She sees $c = 1101$. Could the message be $m' = 0000$? Yes: with the key $k' = m' \oplus c = 1101$. The same holds for \textbf{every} possible message: for each one there is \textbf{exactly one key} that produces $c$. Since all keys are equally likely, Eve has no reason to prefer one message over another. This is the intuition behind the proof.
\end{framedex}

\begin{framedprop}[label={prop:otp}]{}
  The One-Time Pad is perfectly secret.
\end{framedprop}

\begin{proof}
  We use statement (iii) of \Cref{lem:equivalent}, because it only involves the key. Fix any message $m$ and any ciphertext $c$:
  \[
    \Pr_K[\Enc(K, m) = c] \overset{(a)}{=} \Pr_K[K \oplus m = c] \overset{(b)}{=} \Pr_K[K = m \oplus c] \overset{(c)}{=} 2^{-n}
  \]
  \begin{itemize}
    \item \textbf{(a)} This is the definition of encryption in the One-Time Pad.
    \item \textbf{(b)} We ``solve for $K$'' by XOR-ing both sides with $m$: since $m \oplus m = 0^n$, we get $K \oplus m = c \iff K = m \oplus c$. So the event ``encrypting $m$ gives $c$'' is \textbf{exactly} the event ``the key happens to be $m \oplus c$'': there is one and only one key that does it.
    \item \textbf{(c)} $m \oplus c$ is one specific string and the key is uniform, so the probability of picking exactly that key is $\frac{1}{|\Kspace|} = 2^{-n}$.
  \end{itemize}
  The result \textbf{does not depend on $m$}: for any other message $m'$ we get $2^{-n}$ again. So statement (iii) holds and, by the lemma, the One-Time Pad is perfectly secret.
\end{proof}

\begin{framedex}{Intuition with one bit}
  With $n = 1$, ``attack'' $= 0$ and ``retreat'' $= 1$:
  \begin{center}
    \begin{tabular}{lcc}
      \toprule
      \textbf{Message} & \textbf{Key 0} & \textbf{Key 1} \\
      \midrule
      attack (0) & ciphertext 0 & ciphertext 1 \\
      retreat (1) & ciphertext 1 & ciphertext 0 \\
      \bottomrule
    \end{tabular}
  \end{center}
  If Eve sees ciphertext 0, it could be ``attack'' with key 0 or ``retreat'' with key 1, and both keys have probability $\frac{1}{2}$. The ciphertext gives her no reason to prefer one message over the other.
\end{framedex}

A useful byproduct: since $\Pr[C = c] = 2^{-n}$ for every $c$, the ciphertext is \textbf{uniformly random}, whatever the distribution of the message. To Eve it looks like pure noise.

\subsection{Limitations of the One-Time Pad}

If the One-Time Pad is perfect, why don't we use it everywhere? Because perfection has a price:

\begin{table}[H]
  \centering
  \small
  \begin{tabularx}{\textwidth}{@{}>{\raggedright\arraybackslash}p{2.9cm}>{\raggedright\arraybackslash}p{4.2cm}X@{}}
    \toprule
    \textbf{Limitation} & \textbf{Why} & \textbf{Consequence} \\
    \midrule
    Key as long as the message & $|\Kspace| = |\Mspace| = 2^n$ & To send $n$ bits, Alice and Bob must have \textit{already} shared $n$ secret random bits. If they can share the key securely, they could share the message the same way. \\
    One-time only & $c_1 \oplus c_2 = m_1 \oplus m_2$ & Reusing the key reveals the XOR of the two messages; with natural-language text both messages can often be recovered (historical case: the Venona project). \\
    Malleable & $\Dec(k, c \oplus \Delta) = m \oplus \Delta$ & Eve can flip chosen bits of the message without knowing the key: secrecy is not integrity, and this motivates MACs. \\
    Truly uniform key & the proof uses $\Pr[K = k] = 2^{-n}$ & With a biased key, perfect secrecy no longer holds. \\
    \bottomrule
  \end{tabularx}
\end{table}

\section{Shannon's theorem: perfect secrecy requires long keys}

The first limitation is the most serious one. Is it a defect of the One-Time Pad that a smarter scheme could fix? Shannon proved that the answer is \textbf{no}: it holds for \textbf{every} perfectly secret scheme.

\begin{framedthm}[label={thm:shannon}]{Shannon}
  If $\Pi = (\Enc, \Dec)$ is a correct symmetric-key encryption scheme with perfect secrecy, then
  \[
    |\Kspace| \geq |\Mspace|
  \]
  There must be at least as many keys as messages. In terms of lengths: if messages are $n$-bit strings, keys must have \textbf{at least} $n$ bits (equality is allowed, and the One-Time Pad achieves it).
\end{framedthm}

\begin{framedex}{The intuition with three messages and two keys}
  Suppose there are three messages (``attack'', ``retreat'', ``wait'') and only two keys, $k_1$ and $k_2$. Eve intercepts a ciphertext $c$ and, being unbounded, decrypts it with \textbf{every} key: with $k_1$ she gets, say, ``attack'' and with $k_2$ she gets ``retreat''. Since there are only two keys, there are at most two candidate messages, so ``wait'' is \textbf{excluded}: no key decrypts $c$ to ``wait'', so ``wait'' cannot be the message behind $c$.

  Eve has learned something (``it is not \textit{wait}''): before seeing $c$, ``wait'' had some positive probability; after seeing $c$, its probability is zero. Too few keys cannot ``cover'' all messages.
\end{framedex}

\begin{proof}
  By contradiction, suppose that $|\Kspace| < |\Mspace|$.
  \begin{enumerate}
    \item \textbf{Choose the distribution.} Perfect secrecy holds for every distribution, so we choose $M$ uniform over $\Mspace$: every message has probability $\frac{1}{|\Mspace|} > 0$.
    \item \textbf{Fix a ciphertext.} Take any $c$ with $\Pr[C = c] > 0$ (such a $c$ exists, since $C$ always takes some value).
    \item \textbf{Collect the possible decryptions.} Decrypting $c$ with every key gives the set
    \[
      \Mspace(c) := \{\, \Dec(k, c) : k \in \Kspace \,\} \qquad \text{with} \qquad |\Mspace(c)| \leq |\Kspace| < |\Mspace|
    \]
    since every key produces a single message.
    \item \textbf{Find a missing message.} $\Mspace(c)$ is smaller than $\Mspace$, so there is a message $m \in \Mspace \setminus \Mspace(c)$.
    \item \textbf{This message never produces $c$.} By correctness, if $\Enc(k, m) = c$ for some key $k$, then $\Dec(k, c) = m$, that is, $m \in \Mspace(c)$. This is impossible, so no key maps $m$ to $c$.
    \item \textbf{Contradiction.} Then
    \[
      \Pr[M = m \mid C = c] = 0 \neq \frac{1}{|\Mspace|} = \Pr[M = m]
    \]
    which contradicts perfect secrecy. \qedhere
  \end{enumerate}
\end{proof}

Consequences:
\begin{itemize}
  \item \textbf{The One-Time Pad is optimal} $\to$ it has $|\Kspace| = |\Mspace| = 2^n$, the smallest key space allowed by the theorem.
  \item \textbf{Perfect secrecy is inherently impractical} $\to$ to encrypt 1 gigabyte with perfect secrecy, Alice and Bob must have already shared 1 gigabyte of secret, uniformly random key, and they can never reuse it. To protect large amounts of data with short keys (for example 128 bits) we have to weaken the model: an adversary with \textbf{bounded computational power} and \textbf{hardness assumptions}. This is \textbf{computational security}.
\end{itemize}

\section{Going further}\label{sec:going-further}

\subsection{Shannon's characterization theorem}

When the three spaces have the same size, there is a very quick test for perfect secrecy (Katz and Lindell, Theorem 2.10).

\begin{framedthm}{Shannon's characterization}
  If $|\Mspace| = |\Kspace| = |\Cspace|$ and $\Enc$ is deterministic, then $\Pi$ is perfectly secret \textbf{if and only if}:
  \begin{enumerate}
    \item every key is chosen with probability exactly $1/|\Kspace|$;
    \item for every $m \in \Mspace$ and $c \in \Cspace$ there is a \textbf{unique} key $k$ such that $\Enc(k, m) = c$.
  \end{enumerate}
\end{framedthm}

The One-Time Pad passes both conditions immediately: keys are uniform, and $k = m \oplus c$ is the unique key.

\subsection{Exercise: a game-based definition of perfect secrecy}\label{sec:ps-game}

In class this was given as an exercise. The professor's formulation does not choose the bit $b$ at random: it \textbf{fixes} it, and compares two versions of the same game.

\begin{frameddefn}[label={def:ps-game}]{The games $\mathrm{GAME}_{\Pi,\Adv}(b)$}
  For a scheme $\Pi$, an adversary $\Adv$ (even unbounded) and a \textbf{fixed} bit $b \in \{0,1\}$:
  \begin{enumerate}
    \item $\Adv$ outputs two messages $m_0, m_1 \in \Mspace$;
    \item the key is sampled, $k \sample \Kspace$, and $\Adv$ receives $c = \Enc(k, m_b)$;
    \item $\Adv$ outputs a bit $b'$, which is the output of the game.
  \end{enumerate}
  $\Pi$ is \textbf{perfectly secret in the game sense} if for every adversary $\Adv$
  \[
    \Pr\big[\mathrm{GAME}_{\Pi,\Adv}(0) = 1\big] - \Pr\big[\mathrm{GAME}_{\Pi,\Adv}(1) = 1\big] = 0
  \]
  that is, the output of the game has \textbf{the same distribution} whether $m_0$ or $m_1$ is encrypted.
\end{frameddefn}

\begin{framedprob}[label={prob:ps-game}]{}
  Prove that $\Pi$ is perfectly secret (\Cref{def:perfect-secrecy}) if and only if it is perfectly secret in the game sense (\Cref{def:ps-game}).
\end{framedprob}

\begin{proof}[Solution]
  By \Cref{lem:equivalent}, perfect secrecy is equivalent to statement (iii): for all $m, m', c$, $\Pr_K[\Enc(K, m) = c] = \Pr_K[\Enc(K, m') = c]$. So it is enough to show that (iii) is equivalent to the game definition.

  \textbf{(iii) $\Rightarrow$ game.} Fix any adversary $\Adv$ and the pair $m_0, m_1$ it outputs. By (iii), the ciphertext $c = \Enc(K, m_0)$ and the ciphertext $c = \Enc(K, m_1)$ have \textbf{exactly the same distribution}. In both games, $\Adv$ computes $b'$ from $c$ (and from its own random bits, which do not depend on $b$). Applying the same procedure to the same distribution gives the same distribution, so $\Pr[b' = 1]$ is the same in the two games and the difference is $0$. If $\Adv$ chooses $m_0, m_1$ at random, the argument holds for every possible choice, and so also on average.

  \textbf{Game $\Rightarrow$ (iii).} By contradiction, suppose that (iii) fails: there are $m_0, m_1, c^*$ with
  \[
    \Pr_K[\Enc(K, m_0) = c^*] \neq \Pr_K[\Enc(K, m_1) = c^*]
  \]
  Consider the adversary $\Adv$ that outputs exactly these $m_0, m_1$, and outputs $b' = 1$ if and only if the ciphertext it receives is $c^*$. Then
  \[
    \Pr\big[\mathrm{GAME}(0) = 1\big] - \Pr\big[\mathrm{GAME}(1) = 1\big] = \Pr_K[\Enc(K, m_0) = c^*] - \Pr_K[\Enc(K, m_1) = c^*] \neq 0
  \]
  which contradicts the game definition. (This adversary is even efficient: it only compares $c$ with a fixed string. The unbounded power is needed only to \textit{find} $m_0, m_1, c^*$, and the definition quantifies over every adversary, so we may assume it knows them.)
\end{proof}

\paragraph{Link with perfect indistinguishability.} The game definition is also equivalent to perfect indistinguishability (\Cref{def:perfect-indist}), where $b$ is uniform and the adversary must guess it. Write $p_b = \Pr[\mathrm{GAME}(b) = 1]$. When $b$ is uniform:
\[
  \Pr[b' = b] = \frac{1}{2}\Pr[b' = 0 \mid b = 0] + \frac{1}{2}\Pr[b' = 1 \mid b = 1] = \frac{1}{2}(1 - p_0) + \frac{1}{2}p_1 = \frac{1}{2} + \frac{1}{2}(p_1 - p_0)
\]
So $\Pr[b' = b] = \frac{1}{2}$ for every adversary if and only if $p_0 = p_1$ for every adversary. The three formulations (perfect secrecy, the game with fixed $b$, the game with random $b$) say the same thing (see also Katz and Lindell, Lemma 2.6).

This is why the game-based formulation is so important: \textbf{all the computational definitions of the course follow the same pattern}, replacing ``every adversary'' with ``every probabilistic polynomial-time adversary'' and $\frac{1}{2}$ with $\frac{1}{2} + \negl(n)$.

\subsection{The entropy formulation}

The \textit{entropy} $H(X)$ measures how much uncertainty (randomness) a variable contains; the \textit{mutual information} $I(X;Y)$ measures how much information one variable gives about another. Statement (ii) of the lemma says exactly that the mutual information between message and ciphertext is zero, $I(M;C) = 0$. Shannon's theorem also has an \textbf{entropy version}, $H(K) \geq H(M)$:
\[
  H(M) = H(M \mid C) \leq H(M, K \mid C) = H(K \mid C) + \underbrace{H(M \mid K, C)}_{=\,0 \text{ (correctness)}} \leq H(K)
\]
It is stronger than $|\Kspace| \geq |\Mspace|$: it is not enough to have many possible keys, they must also contain at least as much \textit{randomness} as the message.

\subsection{Example and counterexample: the shift cipher}

\begin{framedex}{The shift cipher}
  $\Mspace = \Kspace = \Cspace = \Z_{26}$ (the letters of the alphabet) and $\Enc(k, m) = (m + k) \bmod 26$.
  \begin{itemize}
    \item \textbf{One letter} $\to$ $\Pr_K[\Enc(K, m) = c] = \Pr_K[K = c - m] = 1/26$ for all $m, c$ $\to$ \textbf{perfectly secret}. It is just the One-Time Pad over $\Z_{26}$.
    \item \textbf{Two letters with the same key} $\to$ the ciphertext ``XX'' can come from ``AA'' but \textbf{never} from ``AB'', because the shift preserves the distance between the letters. Statement (iii) fails $\to$ \textbf{not} perfectly secret. And indeed $|\Kspace| = 26 < 26^2 = |\Mspace|$, exactly as Shannon's theorem predicts.
  \end{itemize}
\end{framedex}

% ============================================================
%  Information-theoretic message authentication
% ============================================================

\section{Information-theoretic message authentication}

\paragraph{From secrecy to integrity.} Shannon's theorem tells us that perfect secrecy is possible but expensive: the key must be at least as long as the message. Even if we accept this cost, the One-Time Pad gives us \textbf{confidentiality only}. It is malleable: by flipping bits of the ciphertext, Eve flips the same bits of the message, and Bob cannot notice:
\[
  \Dec(k, c \oplus \Delta) = m \oplus \Delta
\]
So, still \textbf{without any computational assumption}, the natural next question is: can we guarantee \textbf{integrity}, so that Bob always detects a modified message? This leads to \textbf{information-theoretic message authentication codes}. We will see that the story repeats itself: integrity against an unbounded adversary is possible, but the key must grow with the number of messages we want to authenticate.

\subsection{The setting}

In this section we \textbf{do not care about confidentiality}: the message travels in clear. What we want is \textbf{integrity}, that is, Bob must be able to detect whether the message was modified on the way.

Alice and Bob share a secret key $K$, chosen uniformly. Alice sends the message together with a \textbf{tag}, computed with a \textbf{keyed} function:
\[
  \tau = \Tag(K, m), \qquad \Tag : \Kspace \times \Mspace \to \Tspace
\]
where $\Tspace$ is the \textbf{tag space}. Here $\Tag$ is deterministic, and usually it is a keyed hash function, $\Tag(K, m) = h_K(m)$.

\begin{figure}[H]
  \centering
  \begin{tikzpicture}[box/.style={draw, rounded corners=2pt, minimum height=7mm, inner sep=4pt}]
    \node (A) at (0,0) {Alice};
    \node[box] (T) at (2.8,0) {$\Tag(K, m)$};
    \node[box] (V) at (8.2,0) {$\Tag(K, m) \overset{?}{=} \tau$};
    \node (B) at (11,0) {Bob};
    \node (K1) at (2.8,-1.2) {$K$};
    \node (K2) at (8.2,-1.2) {$K$};
    \draw[msg] (A) -- node[above] {$m$} (T);
    \draw[msg] (T) -- node[above] {$(m, \tau)$} (V);
    \draw[msg] (V) -- (B);
    \draw[msg] (K1) -- (T);
    \draw[msg] (K2) -- (V);
  \end{tikzpicture}
  \caption{An integrity-only MAC: the message travels in clear, together with its tag.}
\end{figure}

The protocol, step by step:
\begin{enumerate}
  \item \textbf{Alice} computes the tag of the message, $\tau = \Tag(K, m)$.
  \item Alice sends the \textbf{pair} $(m, \tau)$.
  \item \textbf{Bob} receives $(m, \tau)$, \textbf{recomputes} $\Tag(K, m)$ with the shared key and compares it with $\tau$: if they match he accepts the message, otherwise he rejects it.
\end{enumerate}

\paragraph{The threat.} Eve is an \textbf{active} adversary. She intercepts $(m, \tau)$ and tries to replace it with a pair $(m^*, \tau^*)$, with $m^* \neq m$, that Bob accepts. This is a \textbf{forgery}. The scheme is good if forging is hard \textbf{even after seeing valid pairs}.

\begin{framedobs}{Why the key is essential}
  Suppose we used a hash function \textbf{without a key}, a fixed public function $h$, and set $\tau = h(m)$. Eve intercepts $(m, \tau)$, replaces the message with $m^*$, computes $h(m^*)$ herself (the function is public) and sends $(m^*, h(m^*))$. Bob recomputes $h(m^*)$, sees that it matches, and accepts the forged message.

  With a key, instead, the key \textbf{selects which function} of the family $\{h_k\}$ is used. Eve knows the family (Kerckhoffs's principle) but not which function was chosen, so she cannot compute the correct tag for $m^*$.
\end{framedobs}

Since only the key holders can compute valid tags, a valid tag guarantees both \textbf{integrity} (the message was not modified) and \textbf{authenticity} (it was produced by someone holding the key, that is, by Alice or Bob). It is the shared secret key that gives both properties together.

\subsection{The security definition}

\begin{frameddefn}[label={def:mac-security}]{$t$-time $\varepsilon$-statistical security}
  A MAC $\Pi = (\Tag)$ is $t$-time $\varepsilon$-statistically secure if, for all \textbf{distinct} messages $m, m_1, \ldots, m_t \in \Mspace$ and all tags $\tau, \tau_1, \ldots, \tau_t \in \Tspace$:
  \[
    \Pr_K\big[\Tag(K, m) = \tau \;\big|\; \Tag(K, m_1) = \tau_1, \ldots, \Tag(K, m_t) = \tau_t\big] \leq \varepsilon
  \]
\end{frameddefn}

In words: Eve has seen $t$ valid pairs $(m_i, \tau_i)$. For any \textbf{new} message $m$ and any tag $\tau$ she proposes, the probability that $\tau$ is the correct tag is at most $\varepsilon$.
\begin{itemize}
  \item $t$ $\to$ how many valid pairs the adversary sees (bigger is better);
  \item $\varepsilon$ $\to$ the forging probability (smaller is better);
  \item the adversary is \textbf{unbounded}, which is why the security is called ``statistical''.
\end{itemize}

\begin{framedobs}[label={obs:eps-lower-bound}]{We can never get $\varepsilon = 0$}
  In fact $\varepsilon \geq \frac{1}{|\Tspace|}$ always. Fix everything except $\tau$. The tag of $m$ takes some value, so the probabilities over all possible tags sum to $1$:
  \[
    \sum_{\tau \in \Tspace} \Pr\big[\Tag(K, m) = \tau \mid \ldots\big] = 1
  \]
  A sum of $|\Tspace|$ terms equal to $1$ has at least one term $\geq \frac{1}{|\Tspace|}$, and guessing that tag succeeds with at least that probability. So, to get a small $\varepsilon$, we need long tags.
\end{framedobs}

We often write $\varepsilon = 2^{-\lambda}$, where $\lambda = \log_2 \frac{1}{\varepsilon}$ is the \textbf{security level in bits}: forging is as hard as guessing a random $\lambda$-bit string. For example, $\lambda = 128$ gives $\varepsilon = 2^{-128}$. With this notation, \Cref{obs:eps-lower-bound} says that tags must have \textbf{at least $\lambda$ bits}.

\begin{framedprob}[label={prob:otp-mac}]{The One-Time Pad is not a MAC}
  Let $\Mspace = \Kspace = \Tspace = \{0,1\}^n$ and $\Tag(k, m) = k \oplus m$. Show that this MAC is \textbf{not} $1$-time $\varepsilon$-statistically secure for any $\varepsilon < 1$.
\end{framedprob}

\begin{proof}[Solution]
  Think of the forgery game: Eve sees one valid pair $(m, \tau)$ with $\tau = \Tag(k, m)$, and must output $(m', \tau')$ with $m' \neq m$ and $\tau' = \Tag(k, m')$.
  \begin{enumerate}
    \item \textbf{Recover the key.} From $\tau = k \oplus m$, XOR-ing both sides with $m$ gives $k = \tau \oplus m$. One valid pair reveals the whole key.
    \item \textbf{Forge.} Choose any $m' \neq m$ and output $\tau' = (\tau \oplus m) \oplus m' = k \oplus m'$, which is exactly the correct tag of $m'$.
  \end{enumerate}
  Eve wins with probability $1$. In the language of \Cref{def:mac-security}: given $\Tag(K, m) = \tau$, the key is determined, so
  \[
    \Pr_K\big[\Tag(K, m') = \tau \oplus m \oplus m' \;\big|\; \Tag(K, m) = \tau\big] = 1
  \]
  and the scheme is $1$-time $\varepsilon$-secure only for $\varepsilon = 1$, that is, not secure at all.
\end{proof}

\textbf{The lesson.} The One-Time Pad is perfectly secret, yet it is useless for integrity. Secrecy asks that the output reveals nothing about the \textit{message}; a MAC needs more: after one tag, the tags of \textit{other} messages must still be unpredictable. Here the tag of $m$ and the tag of $m'$ are tied by the fixed relation $\Tag(k, m) \oplus \Tag(k, m') = m \oplus m'$. This is exactly what pairwise independence (below) rules out.

\subsection{The price: long keys (again)}

As with perfect secrecy, statistical security is expensive.

\begin{framedthm}[label={thm:mac-key-size}]{}
  Any $t$-time $2^{-\lambda}$-statistically secure MAC has a key of \textbf{at least} $(t+1)\lambda$ bits.
\end{framedthm}

The proof was omitted in class. The intuition: every valid pair Eve sees can ``use up'' about $\lambda$ bits of the key's uncertainty. To still have $\lambda$ bits of uncertainty about the \textbf{next} tag after $t$ pairs, the key needs $(t+1)\lambda$ bits, that is, $\lambda$ bits for every message to authenticate, plus $\lambda$. The key grows \textbf{linearly with the number of messages}: this is the MAC analogue of Shannon's theorem.

\subsection{Pairwise independent hash functions}

\begin{frameddefn}[label={def:pairwise}]{Pairwise independent hash family}
  A family of functions $\Hfam = \{h_k : \Mspace \to \Tspace\}_{k \in \Kspace}$ is \textbf{pairwise independent} if, for all $m \neq m'$, when $K$ is uniform the pair
  \[
    \big(h_K(m),\, h_K(m')\big) \text{ is uniform over } \Tspace \times \Tspace
  \]
  that is, $\Pr_K[h_K(m) = \tau \wedge h_K(m') = \tau'] = \frac{1}{|\Tspace|^2}$ for all $\tau, \tau'$.
\end{frameddefn}

\textbf{Intuition.} Pick a random function from the family. Its outputs on \textbf{any two different inputs} look like two independent, uniform values: knowing $h_K(m)$ tells you \textbf{nothing} about $h_K(m')$. The independence is between the outputs on two \textit{inputs}, not between two functions.

\textbf{Consequence: each output is uniform.} Summing the joint probability over $\tau'$:
\[
  \Pr_K[h_K(m) = \tau] = \sum_{\tau' \in \Tspace} \frac{1}{|\Tspace|^2} = \frac{1}{|\Tspace|}
\]

\begin{framedthm}[label={thm:pairwise-mac}]{}
  If $\Hfam$ is pairwise independent, then $\Tag(K, m) = h_K(m)$ is a \textbf{1-time $\frac{1}{|\Tspace|}$-statistically secure} MAC.
\end{framedthm}

\begin{proof}
  Eve sees one pair $(m, \tau)$ and tries $(m', \tau')$ with $m' \neq m$. By the definition of conditional probability:
  \[
    \Pr\big[\Tag(K, m') = \tau' \mid \Tag(K, m) = \tau\big]
    = \frac{\Pr[h_K(m') = \tau' \wedge h_K(m) = \tau]}{\Pr[h_K(m) = \tau]}
    = \frac{1/|\Tspace|^2}{1/|\Tspace|} = \frac{1}{|\Tspace|}
  \]
  The numerator uses pairwise independence (the pair is uniform over $\Tspace^2$); the denominator uses the consequence above (each output is uniform). Eve's chance is exactly that of a random guess, which is optimal by \Cref{obs:eps-lower-bound}.
\end{proof}

\subsection{A concrete construction: $am + b \bmod p$}

\begin{framedprop}[label={prop:affine}]{}
  Let $p$ be prime, $\Mspace = \Tspace = \Z_p$, $\Kspace = \Z_p^2$, and
  \[
    h_{(a,b)}(m) = a \cdot m + b \pmod p
  \]
  Then $\Hfam = \{h_{(a,b)}\}_{(a,b) \in \Z_p^2}$ is pairwise independent.
\end{framedprop}

\begin{proof}
  Fix $m \neq m'$ and $\tau, \tau'$. We count how many keys give tag $\tau$ to $m$ \textbf{and} tag $\tau'$ to $m'$. The two conditions form a linear system in the unknowns $a, b$:
  \[
    \begin{pmatrix} m & 1 \\ m' & 1 \end{pmatrix} \begin{pmatrix} a \\ b \end{pmatrix} = \begin{pmatrix} \tau \\ \tau' \end{pmatrix}
  \]
  The determinant is $m - m' \neq 0$. Since $p$ is \textbf{prime}, $\Z_p$ is a field, so every nonzero element is invertible and the matrix is invertible. The system has \textbf{exactly one} solution $(a, b)$, and the key is uniform over $p^2$ possibilities:
  \[
    \Pr_{(a,b)}\big[h(m) = \tau \wedge h(m') = \tau'\big] = \frac{1}{p^2} = \frac{1}{|\Tspace|^2}
  \]
  It is the same ``exactly one key'' argument as for the One-Time Pad.
\end{proof}

\paragraph{Why $p$ must be prime.} To solve the system we must \textbf{divide} by $m - m'$. Modulo a prime, every nonzero number has an inverse. Modulo a composite number this fails: modulo $6$, for example, $2$ has no inverse, and the system could have zero or several solutions.

\paragraph{Geometric intuition.} $y = ax + b$ is the equation of a \textbf{line}. Choosing the key means choosing a \textbf{random line} over $\Z_p$, and the tag of $m$ is the \textbf{height of the line at the point $m$}. Every message-tag pair that Eve sees reveals \textbf{one point} of the line.

\begin{framedobs}{The key can be used only once}
  As with the One-Time Pad, this MAC is secure only if each key authenticates \textbf{a single message}.
  \begin{itemize}
    \item \textbf{One pair} $\to$ one point. Exactly $p$ lines pass through it, one for each slope, and each gives a different tag for any other message $m^*$. So Eve can only guess, with probability $\frac{1}{p}$.
    \item \textbf{Two pairs} $\to$ two points, and only \textbf{one line} passes through two points. Eve solves the system, recovers the key, and can forge a valid tag for any message.
  \end{itemize}
  The key has $2\log p$ bits, that is $(t+1)\lambda$ with $t = 1$ and $\lambda = \log p$: by \Cref{thm:mac-key-size}, it is exactly enough for one message, and no more.
\end{framedobs}

\begin{framedex}{An attack with two pairs}
  With $p = 7$ and secret key $(a, b) = (3, 5)$, Eve sees:
  \begin{itemize}
    \item $m_1 = 2$, with $\tau_1 = 3 \cdot 2 + 5 = 11 \equiv 4$;
    \item $m_2 = 4$, with $\tau_2 = 3 \cdot 4 + 5 = 17 \equiv 3$.
  \end{itemize}
  Subtracting the two equations removes $b$:
  \[
    a = \frac{\tau_1 - \tau_2}{m_1 - m_2} = \frac{4 - 3}{2 - 4} = \frac{1}{-2} \equiv 1 \cdot 5^{-1} \equiv 3 \pmod 7,
    \qquad b = \tau_1 - a m_1 = 4 - 6 \equiv 5 \pmod 7
  \]
  Here $-2 \equiv 5$ and $5^{-1} \equiv 3$ because $5 \cdot 3 = 15 \equiv 1 \pmod 7$. Eve has recovered the whole key $(3, 5)$.

  With a single pair, instead, Eve only knows that $2a + b \equiv 4$: this equation has $7$ solutions, one for each value of $a$, and each gives a different tag for a new message.
\end{framedex}

\subsection{Beyond the lecture: $t$ messages with polynomials}

To authenticate more messages, replace the line with a richer curve: a random \textbf{polynomial of degree $t$},
\[
  h_{(a_0, \ldots, a_t)}(m) = a_0 + a_1 m + a_2 m^2 + \cdots + a_t m^t \pmod p
\]
where the key is the list of the $t+1$ coefficients. The geometry generalizes: through $2$ points passes a single line, through $3$ points a single parabola, and through $t+1$ points a single polynomial of degree $t$. The reason is the Vandermonde matrix.

\begin{frameddefn}{Vandermonde matrix}
  The Vandermonde matrix of the points $x_0, \ldots, x_t$ is the matrix whose $i$-th row contains the powers of $x_i$:
  \[
    V = \begin{pmatrix}
      1 & x_0 & x_0^2 & \cdots & x_0^t \\
      1 & x_1 & x_1^2 & \cdots & x_1^t \\
      \vdots & \vdots & \vdots & & \vdots \\
      1 & x_t & x_t^2 & \cdots & x_t^t
    \end{pmatrix},
    \qquad \det V = \prod_{i < j} (x_j - x_i)
  \]
  The determinant is nonzero if and only if the points are distinct. Hence through $t+1$ distinct points passes \textbf{exactly one} polynomial of degree $t$: its coefficients are the unique solution of $V \cdot (a_0, \ldots, a_t)^\top = (\tau_0, \ldots, \tau_t)^\top$.
\end{frameddefn}

For $t = 1$ the Vandermonde matrix is $\begin{pmatrix} 1 & m \\ 1 & m' \end{pmatrix}$: it is the matrix of \Cref{prop:affine} with the columns swapped, which only changes the sign of the determinant.

Consequences:
\begin{itemize}
  \item after seeing $t$ pairs, $p$ polynomials are still possible and the next tag is uniform $\to$ this is a \textbf{$t$-time $\frac{1}{p}$-secure MAC};
  \item with $t+1$ pairs Eve reconstructs the polynomial, that is, the key;
  \item the key has $(t+1)\log p$ bits, which matches the lower bound $(t+1)\lambda$ of \Cref{thm:mac-key-size} with $\lambda = \log p$: the construction is \textbf{optimal}.
\end{itemize}

% ============================================================
%  Randomness extraction and the Leftover Hash Lemma
% ============================================================

\section{Randomness extraction}

\subsection{The problem: where do uniform keys come from?}

Every scheme we have seen so far, the One-Time Pad, Shannon's theorem and the information-theoretic MACs, starts from the same assumption: \textbf{the key is perfectly uniform}. In the proof of \Cref{prop:otp}, the key step was $\Pr[K = m \oplus c] = 2^{-n}$, which holds only if every key has the same probability; with a biased key, perfect secrecy is lost.

In the real world, though, randomness comes from physical sources:
\begin{itemize}
  \item electronic or thermal noise;
  \item the timing of keyboard and mouse events, or of network packets;
  \item hardware events.
\end{itemize}
These sources are \textbf{unpredictable}, but they are \textbf{not uniform}: some outcomes are more likely than others, and consecutive outputs may be correlated. We cannot use them directly as keys. So the question is:

\begin{center}
  \itshape Given a source that is random but not uniform, can we turn it into bits that are (almost) uniform?
\end{center}

This process is called \textbf{randomness extraction}, and a function that does it is called an \textbf{extractor}. It works like a water filter: a lot of ``dirty'' water goes in, less but clean water comes out.

\begin{figure}[H]
  \centering
  \begin{tikzpicture}[box/.style={draw, rounded corners=2pt, minimum height=8mm, inner sep=5pt}]
    \node[align=center] (X) at (0,0) {$n$ ``dirty'' bits\\[-2pt]{\scriptsize (physical source $X$)}};
    \node[box] (E) at (4,0) {$\Ext$};
    \node[align=center] (Y) at (8,0) {$\ell$ clean bits\\[-2pt]{\scriptsize (almost uniform $\to$ key)}};
    \draw[msg] (X) -- (E);
    \draw[msg] (E) -- (Y);
  \end{tikzpicture}
  \caption{An extractor turns imperfect randomness into (almost) uniform bits.}
\end{figure}

\begin{framedobs}{An extractor does not create randomness}
  The output cannot contain \textbf{more} randomness than the input: an extractor \textbf{concentrates} the randomness already present in the source, it does not multiply it. Expanding a few random bits into many bits that only \textit{look} random is possible only against computationally bounded adversaries, with \textit{pseudorandom generators}, which we will see in the computational part of the course.
\end{framedobs}

\paragraph{The big picture: extract, expand, refresh.} In class the problem was framed as a pipeline with three steps, which is also how real systems (for example the random number generator of an operating system) are organized:
\begin{enumerate}
  \item \textbf{Extract} $\to$ from an imperfect physical source, obtain a few bits that are (almost) uniform. This is the topic of this section.
  \item \textbf{Expand} $\to$ deterministically stretch those few bits into as many bits as needed, which only \textit{look} uniform. This needs a pseudorandom generator (\Cref{def:prg}), and works only against efficient adversaries.
  \item \textbf{Refresh} $\to$ from time to time, mix new extracted randomness into the state, so that an adversary who learns the state at some point cannot predict the future outputs forever.
\end{enumerate}
Security depends on the randomness being both \textbf{uniform} and \textbf{fresh}: extraction takes care of the first, refreshing of the second.

\subsection{A first, simple case: the biased coin}

Let's start from the simplest possible weak source: a \textbf{biased coin} $B$, with $\Pr[B = 1] = p$ for some unknown $0 < p < 1$. We assume that the flips are \textbf{independent} and that the bias $p$ is \textbf{always the same}.

\begin{framedalgo}{Von Neumann's extractor}
  Flip the coin twice and look at the pair $b_1 b_2$ (one \textit{round}):
  \begin{itemize}
    \item $01$ $\to$ output $0$;
    \item $10$ $\to$ output $1$;
    \item $00$ or $11$ $\to$ output nothing: the round \textbf{fails}, throw the pair away and flip twice again.
  \end{itemize}
\end{framedalgo}

The four possible outcomes of a round:

\begin{table}[H]
  \centering
  \small
  \begin{tabular}{llc}
    \toprule
    \textbf{Pair} & \textbf{What happens} & \textbf{Probability} \\
    \midrule
    $01$ & output $0$ & $(1-p)\,p$ \\
    $10$ & output $1$ & $p\,(1-p)$ \\
    $00$ & no output, flip again & $(1-p)^2$ \\
    $11$ & no output, flip again & $p^2$ \\
    \bottomrule
  \end{tabular}
\end{table}

\paragraph{Why the output bit is uniform.} The two useful pairs have exactly the same probability:
\[
  \Pr[01] = (1-p)\,p = p\,(1-p) = \Pr[10]
\]
So, whenever a bit comes out, it is \textbf{exactly} uniform, whatever the bias $p$ is.

\paragraph{Why the procedure terminates.} A round fails with probability $(1-p)^2 + p^2 = 1 - 2p(1-p)$, so $m$ rounds in a row all fail with probability
\[
  \big(1 - 2p(1-p)\big)^m \xrightarrow[m \to \infty]{} 0
\]
which goes to zero exponentially fast. On average we need $\frac{1}{2p(1-p)}$ rounds.

The two facts play \textbf{different roles}: the symmetry $\Pr[01] = \Pr[10]$ makes the bit \textbf{uniform}, while the vanishing failure probability only guarantees that the procedure \textbf{stops}.

\begin{framedex}{A very biased coin}
  With $p = 0.9$: $\Pr[11] = 0.81$ and $\Pr[00] = 0.01$, so a round fails with probability $0.82$, and $50$ rounds in a row fail with probability $0.82^{50} \approx 0.00005$. On average we need $\frac{1}{0.18} \approx 5.6$ rounds. A lot of flips are wasted, but when a bit finally comes out, $0$ and $1$ are equally likely ($0.09$ each).
\end{framedex}

\paragraph{The limit of this approach.} Von Neumann's trick works only because of two strong assumptions: the flips are \textbf{independent}, and they have the \textbf{same bias}. Real sources satisfy neither: their outputs can be correlated, and their distribution can be arbitrary. We need a method that works for \textbf{any} source $X$, without knowing its distribution.

\subsection{How much randomness does a source contain? Min-entropy}

If we want to handle arbitrary sources, we first need to know \textbf{which sources we can hope to extract from}. Some sources clearly contain no randomness at all: if $X$ is constant, there is nothing to extract. So we need a \textbf{measure} of how random a source is.

\paragraph{Which measure?} In cryptography the question that matters is: \textbf{how hard is it for an attacker to guess the value of $X$ at the first attempt?} The best strategy is to bet on the \textbf{most likely value}, which succeeds with probability $\max_x \Pr[X = x]$. The smaller this probability, the more random the source is for the attacker.

A probability like $\frac{1}{256}$ is not very telling, so we compare it with tossing fair coins: guessing $1$ coin succeeds with probability $\frac{1}{2}$, guessing $8$ coins with probability $\frac{1}{2^8} = \frac{1}{256}$, guessing $k$ coins with probability $\frac{1}{2^k}$. Measuring the guessing probability in ``number of coins'' gives the \textbf{min-entropy}.

\begin{frameddefn}[label={def:min-entropy}]{Min-entropy}
  The min-entropy of a random variable $X$ is
  \[
    \Hmin(X) = -\log_2 \max_x \Pr[X = x]
  \]
  Equivalently, $\Hmin(X) \geq k$ if and only if every value of $X$ has probability at most $2^{-k}$.
\end{frameddefn}

So $\Hmin(X) = k$ means that \textbf{guessing $X$ at the first attempt is exactly as hard as guessing $k$ uniformly random bits} (the logarithm just converts the probability into a number of coins). Note that $X$ can be much longer than $k$ bits: a $1000$-bit string whose most likely value has probability $\frac{1}{256}$ is worth, for the attacker, only $8$ random bits. Min-entropy measures how much \textit{real} randomness the source contains, not how long it is.

\begin{table}[H]
  \centering
  \small
  \begin{tabular}{lcc}
    \toprule
    \textbf{Source $X$} & \textbf{Best guess succeeds with} & \textbf{Min-entropy} \\
    \midrule
    one fair coin & $\frac{1}{2}$ & $1$ bit \\
    $8$ fair coins & $\frac{1}{256}$ & $8$ bits \\
    uniform over $\{0,1\}^n$ & $2^{-n}$ & $n$ bits (the maximum) \\
    uniform over a set $S$ & $\frac{1}{|S|}$ & $\log_2 |S|$ bits \\
    a constant & $1$ & $0$ bits \\
    a coin that gives $1$ with probability $0.9$ & $0.9$ & $\approx 0.15$ bits \\
    \bottomrule
  \end{tabular}
\end{table}

\begin{framedex}{Passwords}
  Suppose there are millions of possible passwords, but $1\%$ of the people use ``123456''. An attacker who tries ``123456'' succeeds once in a hundred times:
  \[
    \Hmin = -\log_2 \frac{1}{100} \approx 6.6 \text{ bits}
  \]
  Even though there are millions of possible passwords, against an attacker who guesses the password of a random person they are worth less than $7$ coins. It is the most likely value, not the number of values, that makes a source weak.
\end{framedex}

\paragraph{Why ``min'', and why not Shannon entropy?} Shannon entropy is the \textbf{average} of $\log_2 \frac{1}{\Pr[X = x]}$, while min-entropy takes the \textbf{minimum}, that is, it looks only at the worst case, the most likely value. In general $\Hmin(X) \leq H(X)$, and the difference can be huge. Take $X$ that, with probability $\frac{1}{2}$, is $0^n$, and otherwise is uniform over $\{0,1\}^n$:
\begin{itemize}
  \item Shannon entropy: $H(X) \approx \frac{n}{2}$, which looks like a lot of randomness;
  \item min-entropy: $\Hmin(X) \approx 1$, because an attacker who always guesses $0^n$ is right half of the times.
\end{itemize}
Security is about the attacker's best move, so \textbf{min-entropy is the right measure}.

\subsection{The natural goal, and why it is impossible}

Now that we can measure randomness, we can state the goal precisely. We do not know which source we will get (thermal noise, keyboard timings, \ldots): we only know that it has at least $k$ bits of min-entropy. So:

\begin{center}
  \itshape Find a single function $\Ext$ such that, for \textbf{every} source $X$ with $\Hmin(X) \geq k$, the output $\Ext(X)$ is uniform.
\end{center}

Note that $X$ is a \textit{source}, a random process whose values we do not choose, and $\Ext$ is a \textit{function} applied to whatever value comes out. Unfortunately this goal is impossible, even in the most favourable case: extracting \textbf{a single bit} from a source with almost full min-entropy.

\begin{framedprop}[label={prop:no-deterministic}]{}
  For every function $\Ext : \{0,1\}^n \to \{0,1\}$ there is a source $X$ with $\Hmin(X) \geq n-1$ such that $\Ext(X)$ is \textbf{constant}.
\end{framedprop}

The proof uses the \textbf{pigeonhole principle}: if $N$ objects are distributed into $c$ boxes, at least one box contains at least $\frac{N}{c}$ objects (otherwise there would be fewer than $c \cdot \frac{N}{c} = N$ objects in total).

\begin{proof}
  \begin{enumerate}
    \item $\Ext$ splits the $2^n$ inputs into two groups: those mapped to $0$ and those mapped to $1$.
    \item By the pigeonhole principle ($2^n$ objects, $2$ boxes), one of the two groups contains at least $2^{n-1}$ inputs. Call $b$ the output of that group.
    \item Let $X$ be uniform over that group. Every value has probability at most $\frac{1}{2^{n-1}}$, so
    \[
      \Hmin(X) = -\log_2 \max_x \Pr[X = x] \geq -\log_2 \frac{1}{2^{n-1}} = n - 1
    \]
    \item But every value of $X$ is mapped to $b$, so $\Ext(X) = b$ \textbf{always}. \qedhere
  \end{enumerate}
\end{proof}

\begin{framedex}{The parity extractor with $n = 3$}
  Take $\Ext(x_1 x_2 x_3) = x_1 \oplus x_2 \oplus x_3$ (the parity), which seems a good extractor since it mixes all the bits. The $8$ inputs split into two groups:
  \begin{center}
    \begin{tabular}{cc}
      \toprule
      \textbf{Output 0 (even number of 1s)} & \textbf{Output 1 (odd number of 1s)} \\
      \midrule
      $000,\ 011,\ 101,\ 110$ & $001,\ 010,\ 100,\ 111$ \\
      \bottomrule
    \end{tabular}
  \end{center}
  Let $X$ be uniform over $\{000, 011, 101, 110\}$: every value has probability $\frac{1}{4}$, so $\Hmin(X) = 2 = n - 1$. Guessing $X$ is as hard as guessing $2$ coins, and yet $\Ext(X) = 0$ always: we extracted no randomness at all.
\end{framedex}

\begin{framedobs}{Moral}
  The problem is not the source, which is very random, but the fact that the extractor is \textbf{fixed in advance}: once $\Ext$ is known, an adversarial source can always be built against it, by concentrating on inputs that $\Ext$ maps to the same output. To avoid this, we relax the goal to an output that is only \textbf{close} to uniform (\Cref{sec:sd}), and we let the extractor use a \textbf{random seed}, chosen independently of the source, so that the function is picked at random \textbf{after} the source is fixed (\Cref{sec:seeded}).
\end{framedobs}

Why should a source be ``built against'' the extractor? Two reasons. First, we want a \textbf{guarantee} for every source with enough min-entropy, because we do not know the exact distribution of ours; the proposition shows that a fixed extractor cannot give this guarantee. Second, the source may really be influenced by an adversary: a server that uses the arrival times of network packets as randomness receives packets from the attacker too, who can choose when to send them. The attacker does not control the source completely, so it remains unpredictable, but she can \textit{shape} its distribution.

\paragraph{Three ways out.} In class three ways around \Cref{prop:no-deterministic} were listed:
\begin{enumerate}
  \item \textbf{Several independent sources} $\to$ extract from two (or more) sources $X_1, X_2$, each with enough min-entropy, that are \textbf{independent} of each other. A single adversarial source can be built against a fixed function, but two independent ones cannot ``coordinate''. These are called \textit{two-source} (or multi-source) extractors; we do not study them.
  \item \textbf{Cryptographic hash functions} $\to$ apply a hash function such as SHA-256 to the source. It works well in practice, but proving it requires \textbf{strong assumptions} on the hash function (it is modelled as a truly random function), so it is not an information-theoretic solution.
  \item \textbf{A random seed} $\to$ give the extractor a short uniform seed, independent of the source. This is the approach we follow in the rest of the section (\Cref{sec:seeded}).
\end{enumerate}

\subsection{Measuring ``close to uniform'': statistical distance}\label{sec:sd}

The first lesson of \Cref{prop:no-deterministic} is that exact uniformity is too much to ask. To relax it to ``almost uniform'' we need a way to measure how far apart two distributions are.

\begin{frameddefn}[label={def:sd}]{Statistical distance}
  The statistical distance between two random variables $X$ and $X'$ over the same set is
  \[
    \SD(X; X') = \frac{1}{2} \sum_x \big|\Pr[X = x] - \Pr[X' = x]\big|
  \]
  We write $X \approx_\varepsilon X'$, read ``$\varepsilon$-close'', when $\SD(X; X') \leq \varepsilon$.
\end{frameddefn}

The meaning is what matters: the statistical distance is exactly the \textbf{best advantage of an unbounded adversary} in telling the two distributions apart,
\[
  \SD(X; X') = \max_{\Adv} \Big|\Pr[\Adv(x) = 1 : x \leftarrow X] - \Pr[\Adv(x) = 1 : x \leftarrow X']\Big|
\]
So ``$\varepsilon$-close to uniform'' means that no adversary, however powerful, can distinguish the output from truly uniform bits with advantage more than $\varepsilon$. It is the same ``quiz'' idea of the security definitions.

Useful properties:
\begin{itemize}
  \item $0 \leq \SD \leq 1$;
  \item triangle inequality: $\SD(X; Z) \leq \SD(X; Y) + \SD(Y; Z)$;
  \item \textbf{processing cannot help}: $\SD(f(X); f(X')) \leq \SD(X; X')$ for every function $f$.
\end{itemize}

\subsection{Adding a seed: seeded extractors}\label{sec:seeded}

The second lesson was that the extractor needs randomness that the source does not know. So we give it a short, \textbf{truly uniform seed} $S$, chosen \textbf{independently of the source}:
\[
  \Ext : \underbrace{\{0,1\}^d}_{\text{seed}} \times \underbrace{\{0,1\}^n}_{\text{source}} \to \underbrace{\{0,1\}^\ell}_{\text{output}}
\]

\paragraph{A trap to avoid.} If we only asked ``$\Ext(S, X)$ is close to uniform'', there would be a useless solution: $\Ext(S, X) = S$, which ignores the source and copies the seed. It is uniform, but it extracts nothing from $X$. To rule this out, we ask the output to look uniform \textbf{even to someone who sees the seed}. As a bonus, the seed can then be \textbf{public}.

\begin{frameddefn}[label={def:extractor}]{Seeded extractor}
  $\Ext : \{0,1\}^d \times \{0,1\}^n \to \{0,1\}^\ell$ is a \textbf{$(k, \varepsilon)$-extractor} if, for every source $X$ with $\Hmin(X) \geq k$:
  \[
    \big(S,\ \Ext(S, X)\big) \approx_\varepsilon \big(S,\ U_\ell\big)
  \]
  where $S$ is uniform over $\{0,1\}^d$ and independent of $X$, and $U_\ell$ is uniform over $\{0,1\}^\ell$.
\end{frameddefn}

In words: even knowing the seed, the output is indistinguishable from fresh uniform bits, except with advantage $\varepsilon$. (This notion is also called a \textit{strong} seeded extractor.)

\begin{framedex}{Alice, Bob and Eve}
  Alice and Bob want to use the One-Time Pad to exchange one bit (``attack'' or ``retreat''), so they need a \textbf{uniform 1-bit key}. Alice's computer generates it from the last $3$ bits of the arrival time of a network packet: a source $X$ of $3$ bits. Eve cannot read Alice's computer, but she can \textbf{send packets} to it and partly influence their arrival times.

  \textbf{Scenario 1: fixed extractor.} Alice uses the parity, and Eve knows it (Kerckhoffs's principle). Eve times her packets so that the arrival times fall on $\{000, 011, 101, 110\}$, at random among them. The source still has $\Hmin(X) = 2$, so Alice notices nothing, but the key is always $0$. The next day Alice encrypts $c = m \oplus 0 = m$: the ciphertext is the message, and Eve reads it in clear.

  \textbf{Scenario 2: seeded extractor.} Alice uses $h_s(x) = (s_1 x_1) \oplus (s_2 x_2) \oplus (s_3 x_3)$, the XOR of the bits selected by a seed $s$ chosen at random \textbf{after} Eve has prepared her packets. Against the same source:
  \begin{center}
    \small
    \begin{tabular}{cccc}
      \toprule
      \textbf{Seed $s$} & \textbf{Bits used} & \textbf{Output on $000, 011, 101, 110$} & \textbf{Key} \\
      \midrule
      $100$ & $x_1$ & $0,0,1,1$ & uniform \\
      $010$ & $x_2$ & $0,1,0,1$ & uniform \\
      $001$ & $x_3$ & $0,1,1,0$ & uniform \\
      $110$ & $x_1 \oplus x_2$ & $0,1,1,0$ & uniform \\
      $101$ & $x_1 \oplus x_3$ & $0,1,0,1$ & uniform \\
      $011$ & $x_2 \oplus x_3$ & $0,0,1,1$ & uniform \\
      $111$ & parity & $0,0,0,0$ & constant \\
      $000$ & none & $0,0,0,0$ & constant \\
      \bottomrule
    \end{tabular}
  \end{center}
  With $6$ seeds out of $8$ the key is perfectly uniform: Eve's trap was built against \textit{one} function, and the function is drawn at random. With realistic sizes (hundreds of bits), the Leftover Hash Lemma below guarantees that the fraction of ``unlucky'' seeds becomes negligible.

  The seed may even be public, but it must be chosen \textbf{independently of the source}: if Alice always used the same seed, known in advance, we would be back to Scenario 1.
\end{framedex}

\subsection{Random hash functions and collision probability}

Do seeded extractors exist? A natural candidate comes from the MAC section: \textbf{pairwise independent hash functions}. The seed $s$ selects a function from the family, and
\[
  \Ext(s, x) = h_s(x)
\]
\textbf{Intuition.} A random function from a pairwise independent family \textbf{spreads its inputs evenly}: two different inputs land on the same output only with probability $2^{-\ell}$. A source with high min-entropy has many inputs, none too likely, so after hashing they get spread almost uniformly.

To prove it, we need to bound the statistical distance of $(S, h_S(X))$ from uniform. Computing it directly is hard, because of the absolute values. There is a quantity that is much easier to compute for hash functions and that still controls the distance: the \textbf{collision probability}.

\begin{frameddefn}{Collision probability}
  Let $Y'$ be an independent copy of $Y$. The collision probability of $Y$ is
  \[
    \Col(Y) = \Pr[Y = Y'] = \sum_y \Pr[Y = y]^2
  \]
  that is, the probability that two independent samples of $Y$ are equal.
\end{frameddefn}

\textbf{Why it is a good measure.} It is the birthday paradox: two people in a room share their birthday with probability $\frac{1}{365}$ if birthdays are uniform, but much more often if almost everyone was born in January. \textbf{The more unbalanced a distribution is, the more it collides.} Let $N = |\mathcal{Y}|$ be the number of possible values:

\begin{table}[H]
  \centering
  \small
  \begin{tabular}{@{}lll@{}}
    \toprule
    \textbf{Distribution} & $\Col$ & \\
    \midrule
    uniform over $N$ values & $N \cdot \left(\frac{1}{N}\right)^2 = \frac{1}{N}$ & the minimum possible \\
    fair coin & $\frac{1}{4} + \frac{1}{4} = 0.5$ & \\
    coin with bias $90\%$ & $0.81 + 0.01 = 0.82$ & higher \\
    constant & $1$ & the maximum possible \\
    \bottomrule
  \end{tabular}
\end{table}

\subsubsection{Collisions in excess}

The uniform distribution has the \textbf{fewest possible collisions}, and $\Col(Y) = \frac{1}{N}$ if and only if $Y$ is uniform. So the collision probability of any $Y$ over $N$ values can be written as ``uniform plus an extra'':
\[
  \Col(Y) = \frac{1}{N}\,(1 + \delta) \qquad \text{with } \delta \geq 0
\]
\begin{itemize}
  \item \textbf{The $1$} is the part that every distribution has, the collisions of the uniform distribution.
  \item \textbf{The excess $\delta$} measures how much more $Y$ collides than the uniform distribution, that is, \textbf{how far $Y$ is from uniform}: $\delta = 0$ if and only if $Y$ is uniform, and $\delta = N - 1$ when $Y$ is constant.
\end{itemize}

Writing $q_y = \Pr[Y = y] - \frac{1}{N}$ for the \textbf{deviation} of the value $y$ from uniform, we will see in the proof below that
\[
  \delta = N \cdot \sum_y q_y^2
\]
so $\delta$ and the statistical distance measure the \textbf{same thing in two different ways}:
\begin{itemize}
  \item \textbf{the statistical distance} $\SD(Y; U) = \frac{1}{2}\sum_y |q_y|$ uses the \textbf{absolute values} of the deviations: it has a direct meaning (the best advantage of a distinguisher), but it is hard to compute;
  \item \textbf{the excess} $\delta$ uses the \textbf{squares} of the deviations: it is easy to compute for hash functions.
\end{itemize}
The next lemma is the bridge between the two.

\begin{framedlem}[label={lem:col-sd}]{From collisions to statistical distance}
  Let $Y$ take values in a set $\mathcal{Y}$ of size $N$, let $U$ be uniform over $\mathcal{Y}$, and write $\Col(Y) = \frac{1}{N}(1 + \delta)$. Then
  \[
    \SD(Y; U) \leq \frac{1}{2}\sqrt{\delta}
  \]
  In particular, if $\Col(Y) \leq \frac{1}{N}(1 + 4\varepsilon^2)$, then $\SD(Y; U) \leq \varepsilon$.
\end{framedlem}

\begin{proof}
  \leavevmode
  \begin{enumerate}
    \item \textbf{Rewrite the distance with the deviations.} Since $\Pr[U = y] = \frac{1}{N}$ for every $y$:
    \[
      \SD(Y; U) = \frac{1}{2}\sum_y \Big|\Pr[Y = y] - \frac{1}{N}\Big| = \frac{1}{2}\sum_y |q_y|
    \]
    The absolute values are what makes this hard to compute, so we want to move to squares.
    \item \textbf{Cauchy--Schwarz.} A sum of $N$ numbers is at most $\sqrt{N}$ times the square root of the sum of their squares (with equality when the numbers are all equal):
    \[
      \sum_y |q_y| = \sum_y |q_y| \cdot 1 \leq \sqrt{\sum_y q_y^2} \cdot \sqrt{\sum_y 1^2} = \sqrt{N \cdot \sum_y q_y^2}
    \]
    so
    \[
      \SD(Y; U) \leq \frac{1}{2}\sqrt{N \cdot \sum_y q_y^2}
    \]
    \item \textbf{The sum of squares is the excess.} Expand $(a - b)^2 = a^2 - 2ab + b^2$:
    \[
      \sum_y q_y^2 = \underbrace{\sum_y \Pr[Y = y]^2}_{=\,\Col(Y)} - \frac{2}{N}\underbrace{\sum_y \Pr[Y = y]}_{=\,1} + \underbrace{\sum_y \frac{1}{N^2}}_{=\,\frac{1}{N}} = \Col(Y) - \frac{1}{N}
    \]
    and with $\Col(Y) = \frac{1}{N}(1 + \delta)$:
    \[
      \sum_y q_y^2 = \frac{1}{N} + \frac{\delta}{N} - \frac{1}{N} = \frac{\delta}{N}
    \]
    \item \textbf{Substitute:}
    \[
      \SD(Y; U) \leq \frac{1}{2}\sqrt{N \cdot \frac{\delta}{N}} = \frac{1}{2}\sqrt{\delta}
    \]
    \item \textbf{Choose the threshold.} If $\delta \leq 4\varepsilon^2$, then $\SD(Y; U) \leq \frac{1}{2}\sqrt{4\varepsilon^2} = \varepsilon$. \qedhere
  \end{enumerate}
\end{proof}

\begin{framedobs}{Where $1 + 4\varepsilon^2$ comes from}
  The threshold is not the result of a computation: it is chosen \textbf{backwards}, starting from the goal $\SD \leq \varepsilon$.
  \begin{itemize}
    \item \textbf{The $1$} $\to$ the collisions that even the uniform distribution has.
    \item \textbf{The $4\varepsilon^2$} $\to$ the largest excess of collisions compatible with $\SD \leq \varepsilon$, because $\frac{1}{2}\sqrt{\delta} \leq \varepsilon \iff \delta \leq 4\varepsilon^2$. The $4$ comes from the $\frac{1}{2}$ in the definition of the statistical distance: squaring $2\varepsilon$ gives $4\varepsilon^2$.
  \end{itemize}
\end{framedobs}

\begin{framedex}{A coin}
  For a coin, $N = 2$, so $\Col(Y) = \frac{1}{2}(1 + \delta)$ gives $\delta = 2\,\Col(Y) - 1$.
  \begin{center}
    \small
    \begin{tabular}{@{}lcccc@{}}
      \toprule
      \textbf{Coin} & $\Col(Y)$ & $\delta$ & \textbf{Bound} $\frac{1}{2}\sqrt{\delta}$ & \textbf{True} $\SD$ \\
      \midrule
      fair ($50/50$) & $0.5$ & $0$ & $0$ & $0$ \\
      biased ($90/10$) & $0.82$ & $0.64$ & $0.4$ & $\frac{1}{2}(0.4 + 0.4) = 0.4$ \\
      constant ($100/0$) & $1$ & $1$ & $0.5$ & $\frac{1}{2}(0.5 + 0.5) = 0.5$ \\
      \bottomrule
    \end{tabular}
  \end{center}
  With only two values the bound is even exact. With more values it is in general only an upper bound, but that is all we need.
\end{framedex}

So, to prove that hashing gives an extractor, it is enough to show that $(S, h_S(X))$ has \textbf{few collisions in excess}.

\subsection{The Leftover Hash Lemma}

We now have all the pieces: random hashing as the candidate extractor, and the collision probability as the tool to analyse it.

\begin{framedthm}[label={thm:lhl}]{Leftover Hash Lemma}
  Let $\Hfam = \{h_s : \{0,1\}^n \to \{0,1\}^\ell\}_{s \in \{0,1\}^d}$ be a pairwise independent family. If
  \[
    k \geq \ell + 2\log_2\frac{1}{\varepsilon} - 2
  \]
  then $\Ext(s, x) = h_s(x)$ is a $(k, \varepsilon)$-extractor.
\end{framedthm}

\paragraph{The statement used in class.} The professor writes the output length as $m$ and states the lemma with
\[
  k = m + 2\log_2\frac{1}{\varepsilon}
\]
This is the form to remember for the exam. It is the same theorem: since $m + 2\log_2\frac{1}{\varepsilon} \geq m + 2\log_2\frac{1}{\varepsilon} - 2$, the professor's condition implies ours. The version above only gains $2$ bits, by keeping track of the constant $4$ in the lemma.

\textbf{In words.} From a source with $k$ bits of min-entropy, hashing with a random public function gives about
\[
  \ell \approx k - 2\log_2\frac{1}{\varepsilon}
\]
bits that are $\varepsilon$-close to uniform. The $2\log_2\frac{1}{\varepsilon}$ bits we lose are the \textbf{price} of being $\varepsilon$-close instead of exactly uniform.

\begin{proof}
  Let $Y = (S, h_S(X))$, and let $Y' = (S', h_{S'}(X'))$ be an independent copy, where $S'$ and $X'$ are independent copies of $S$ and $X$. We bound $\Col(Y)$ and then apply \Cref{lem:col-sd}.
  \begin{enumerate}
    \item \textbf{Separate the seeds.} The two samples collide only if the seeds are equal \textbf{and} the outputs are equal:
    \[
      \Col(Y) = \Pr[S = S' \wedge h_S(X) = h_{S'}(X')] = \underbrace{\Pr[S = S']}_{=\,2^{-d}} \cdot \Pr[h_S(X) = h_S(X')]
    \]
    Once $S = S'$, the seed is still uniform and independent of $X$ and $X'$.
    \item \textbf{Two ways for the outputs to collide.} Either the inputs themselves are equal, or they are different and the hash maps them to the same output:
    \[
      \Pr[h_S(X) = h_S(X')] \leq \Pr[X = X'] + \Pr[X \neq X' \wedge h_S(X) = h_S(X')]
    \]
    \item \textbf{Equal inputs: this is where min-entropy enters.}
    \[
      \Pr[X = X'] = \sum_x \Pr[X = x]^2 \leq \max_x \Pr[X = x] \cdot \sum_x \Pr[X = x] = 2^{-\Hmin(X)} \leq 2^{-k}
    \]
    \item \textbf{Different inputs: this is where pairwise independence enters.} For fixed $x \neq x'$, the pair $(h_S(x), h_S(x'))$ is uniform over $\{0,1\}^\ell \times \{0,1\}^\ell$, so
    \[
      \Pr_S[h_S(x) = h_S(x')] = \sum_{\tau} \Pr[h_S(x) = \tau \wedge h_S(x') = \tau] = 2^\ell \cdot 2^{-2\ell} = 2^{-\ell}
    \]
    Averaging over $X$ and $X'$, the second term is at most $2^{-\ell}$.
    \item \textbf{Put the pieces together.}
    \[
      \Col(Y) \leq 2^{-d}\big(2^{-k} + 2^{-\ell}\big) = \frac{1}{2^{d+\ell}}\big(1 + 2^{\ell - k}\big)
    \]
    This has exactly the form ``uniform plus an excess'', with $N = 2^{d+\ell}$:
    \begin{itemize}
      \item \textbf{the $1$} comes from the \textbf{different} inputs, which thanks to the hash collide as if they were uniform;
      \item \textbf{the excess $\delta = 2^{\ell - k}$} comes from the \textbf{equal} inputs, that is, from the source repeating the same value. The higher the min-entropy $k$, the smaller it is.
    \end{itemize}
    \item \textbf{Use the hypothesis on $k$.} The lemma needs $\delta \leq 4\varepsilon^2$:
    \[
      2^{\ell - k} \leq 4\varepsilon^2 \iff \ell - k \leq 2 - 2\log_2\frac{1}{\varepsilon} \iff k \geq \ell + 2\log_2\frac{1}{\varepsilon} - 2
    \]
    which is exactly the hypothesis of the theorem (the $-2$ is $\log_2 4$, the same $4$ as in $4\varepsilon^2$). Hence
    \[
      \Col(Y) \leq \frac{1 + 4\varepsilon^2}{2^{d+\ell}}
    \]
    \item \textbf{Conclude with the lemma.} $Y$ takes values in a set of size $2^{d+\ell}$, and the uniform distribution on that set is exactly $(S, U_\ell)$. By \Cref{lem:col-sd}, $\SD\big((S, h_S(X)); (S, U_\ell)\big) \leq \varepsilon$. \qedhere
  \end{enumerate}
\end{proof}

\begin{framedobs}{}
  \begin{itemize}
    \item \textbf{The two hypotheses meet in the proof}: min-entropy controls the collisions of the \textit{source} (step 3), pairwise independence controls the collisions of the \textit{hash} (step 4).
    \item \textbf{Only collisions are needed}: the proof uses pairwise independence only through $\Pr[h_S(x) = h_S(x')] \leq 2^{-\ell}$ for $x \neq x'$. Families with this weaker property are called \textit{universal}, and they are enough. The family of the Alice and Bob example is universal.
    \item \textbf{The same tool twice}: the family $am + b \bmod p$ used for the MACs works here too.
    \item \textbf{A numerical example}: from a source with $k = 256$ bits of min-entropy and a target $\varepsilon = 2^{-64}$, we extract $\ell = 256 - 128 + 2 = 130$ almost-uniform bits.
    \item \textbf{The loss cannot be avoided}: losing about $2\log_2\frac{1}{\varepsilon}$ bits is optimal in general (a lower bound due to Radhakrishnan and Ta-Shma).
    \item \textbf{Where it is used}: deriving keys from Diffie--Hellman shared secrets, which have high min-entropy but are not uniform; privacy amplification in quantum key distribution; extracting keys from physical noise.
  \end{itemize}
\end{framedobs}

\subsection{The thread of the argument}

\begin{table}[H]
  \centering
  \small
  \begin{tabularx}{\textwidth}{@{}cXX@{}}
    \toprule
    \textbf{Step} & \textbf{Problem we face} & \textbf{What it leads to} \\
    \midrule
    1 & keys must be uniform, but real sources are not & extractors \\
    2 & simplest weak source: a biased coin & Von Neumann's extractor (independent flips with the same bias only) \\
    3 & we want any source: from which ones can we extract? & min-entropy, the attacker's best guessing probability \\
    4 & one function that makes every high-min-entropy source uniform & impossible: every fixed function can be fooled \\
    5 & exact uniformity is too much & statistical distance \\
    6 & a fixed function can be fooled & a public random seed: seeded extractors \\
    7 & which function should we use? & random hash functions from a pairwise independent family \\
    8 & the statistical distance is hard to compute directly & collision probability \\
    9 & put everything together & Leftover Hash Lemma: $\ell \approx k - 2\log_2\frac{1}{\varepsilon}$ \\
    \bottomrule
  \end{tabularx}
\end{table}

